\documentclass[acmsmall, screen]{acmart}
\let\Bbbk\relax

\usepackage{quiver}
\usepackage{mathpartir}
% \usepackage{tikz-cd}
\usepackage{enumitem}
\usepackage{wrapfig}
\usepackage{fancyvrb}
\usepackage{comment}
\usepackage{array}

% \theoremstyle{remark} 
% \newtheorem{theorem}{Theorem}
% \newtheorem{corollary}{Corollary}
\theoremstyle{remark}
\newtheorem{remark}{Remark}
\newenvironment{DIFnomarkup}{}{}
% \pagestyle{empty}
% \settopmatter{printfolios=false}

\newcommand{\bibliographypath}{../../paper-new/references}
\newcommand{\defspath}{../../paper-new/defs}


\input{\defspath}


\begin{document}

\title{Denotational Semantics for Effectful Gradually Typed Languages using Guarded Type Theory}
\author{Eric Giovannini}
\affiliation{
  \department{Electrical Engineering and Computer Science}
  \institution{University of Michigan}
  \country{USA}
}
\email{ericgio@umich.edu}
\orcid{0009-0003-6871-1714}


\author{Max S. New}
\affiliation{
  \department{Electrical Engineering and Computer Science}
  \institution{University of Michigan}
  \country{USA}
}
\email{maxsnew@umich.edu}
\orcid{0000-0001-8141-195X}

\begin{abstract}
    Gradually-typed langugaes allow for soundly mixing static and dynamically
    typed programming styles and support migration from a dynamic to static
    typing style. Establishing the desired metatheoretic properties for
    real-world gradually-typed languages has proved challenging.
    %
    In recent work, Giovannini et al. gave a denotational semantics for a simple
    gradually-typed lambda calculus using the tools of guarded type theory, and
    furthermore showed that their semantics was sufficient for establishing the
    crucial metatheoretic property known as \emph{graduality}, which states that
    migrating a program by making the types more precise does not fundamentally
    chance the meaning of the program except to possibly catch more type errors.
    %
    
    In this work, we address two limitations of the work of Giovannini et al.
    First, we generalize their model to allow for languages with algebraic
    effects, such as I/O, global state, nondeterminism, etc. Second, we relate
    the denotational semantics to a more familiar operational semantics.
\end{abstract}

\maketitle

\newif\ifdraft
\drafttrue
\renewcommand{\max}[1]{\ifdraft{\color{blue}[{\bf Max}: #1]}\fi}
\newcommand{\eric}[1]{\ifdraft{\color{orange}[{\bf Eric}: #1]}\fi}
\newcommand{\tingting}[1]{\ifdraft{\color{red}[{\bf Tingting}: #1]}\fi}


\section{Introduction}

\emph{Gradually typed languages} \cite{siek-taha06, tobin-hochstadt06} allow for
safe interoperability between static and dynamic typing disciplines in the same
codebase, and they support the smooth, \emph{gradual} migration from dynamic
typing to static typing via the addition of explicit typing information.
%
% As a result, the programmer can begin with fully dynamically-typed code and
% \emph{gradually} add explicit typing information to migrate portions of the
% codebase to a statically typed style without needing to rewrite the entire
% project in a completely different language.
%
One of the fundamental theorems for gradually typed languages is
\emph{graduality}, also known as the \emph{dynamic gradual guarantee},
originally defined by Siek, Vitousek, Cimini, and Boyland
\cite{siek_et_al:LIPIcs:2015:5031, new-ahmed2018}.
%
Informally, graduality says that migrating code from dynamic to static typing
should only allow for the introduction of static or dynamic type errors, and not
otherwise change the behavior of the program.
%

Establishing graduality for a given cast calculus has traditionally been a
tedious endeavor. One common approach involves defining a step-indexed logical
relation over an operational semantics \cite{TODO}; these methods are inherently
dependent on the particular calculus being modeled and offer little opportunity
for reuse of mathematical abstractions across different developments. Another
method, developed by New and Licata \cite{new-licata-18}, constructs a
denotational semantics using \emph{classical} domain theory. While
compositional, this approach suffers from a lack of expressiveness: classical
domain-theoretic techniques appear incapable of modeling viciously
self-recursive structures such as higher-order store
\cite{Birkedal-Stovring-Thamsborg-2009}, which are common in real-world
gradually-typed languages.
%
Thus, it remained an open question whether one could give a semantics for
gradual typing that combined the expressiveness of the step-indexed approach
with the compositional nature of denotational semantics.

% Graduality has remained a difficult property to establish, with existing
% methods being largely tied to syntactic or operational specifics of the
% language being modeled and offering little opportunity for reuse of common
% mathematical abstractions.

% Recently, Giovannini et al. \cite{gtt-sgdt} combined the expressiveness of the
% step-indexed approach with the compositional nature of denotational semantics
% by...

Recently, Giovannini et al. \cite{gtt-sgdt} answered this question affirmatively
by constructing a denotational semantics for a simple gradually-typed
programming language using the tools of guarded type theory \cite{first-steps},
and proved that their model validates the graduality property. Their model
combines the expressive power of step-indexing techniques with the compositional
nature of denotational semantics.
%
However, the semantics developed by Giovannini et al. suffers from the
limitation that it is a single model for a quite bare-bones gradually-typed
language. Real-world languages can include effects, polymorphism, dependent
types, etc.
%
Furthermore, it is unclear how to relate their model to a more typical
operational semantics for a cast calculus. It is desirable to establish an
adequacy result relating their denotational semantics to an operational model.
More specifically, writing $\sem{M}$ for the denotation of a term $M$, adequacy
states that an equality of denotations $\sem{M} = \sem{N}$ implies that $M$ and
$N$ behave the same operationally. This allows one to use facts proved in the
denotational semantics to reason about the operational behavior of terms.

% In this work, we seek to address these limitations by
In this work, we seek to address the limitations of the prior model by extending
their denotational semantics to account for arbitrary algebraic effects, and
establishing an adequacy theorem relating the resulting denotational model to an
operational semantics.

The level of generality needed to model arbitrary algebraic effects means that
the semantics we produce is correspondingly generic. This is true both
denotationally and operationally. On the denotational side, our guarded
semantics models effectful computations as a kind of \emph{inductive-guarded
interaction tree}. These are a generalization of guarded interaction trees
\cite{guarded-interaction-trees} that allow for inductive occurrences of the
variable in addition to guarded-recursive ones. The (guarded) graduality theorem
we obtain is essentially a statement relating term precision to the behavior of
the interaction trees that the terms denote.
%
Likewise, we construct a generic effectful operational semantics following the
methodology of Kavvos \cite{TODO}. While this approach is applicable to any
signature of effects, what is missing is a way to turn the semantics into a more
traditional ``runnable'' operational semantics. For example, a big-step
semantics for a language with global state would normally involve a transition
relation between pairs of a term and state.

The missing step in both the operational and denotational semantics is supplied
by a suitable notion of algebra that interprets the guarded interaction trees as
a concrete element of an ``observation'' type.

More concretely, the organization of this paper is as follows:
\begin{enumerate}
    \item Given a signature of algebraic effects, we define a cast calculus
    supporting those effects. We give an inequational theory in the form of
    axioms for type and term precision, and define an operational semantics for
    the cast calculus, drawing inspiration from recent work of Kavvos
    \cite{TODO}.
    \item We show how to construct a denotational semantics for such a cast
    calculus by adapting the construction of Giovannini et. al \cite{gtt-sgdt}
    to account for arbitrary algebraic effects.
    \item Following the technique of the aforementioned work by Kavvos, we prove
    an adequacy result relating the operational and denotational models in the
    expected manner.
\end{enumerate}


\subsection{Metatheoretic Preliminaries}

% Working informally in a guarded type theory with ticks and clocks.


\section{A Generic Syntax for Gradual Typing with Effects}\label{sec:syntax}

To begin, we fix a framework for specifying gradual cast calculi with algebraic
effects. We use a syntax based on fine-grain call-by-value presented in Figure
\ref{fig:gtlc-syntax}.

% TODO should we use CBPV syntax instead?
% TODO what about coproducts and lazy products?

% TODO if we are doing fine-grain CBV, do we need return and bind?

\begin{figure}
  \begin{mathpar}
    \begin{array}{rcl}
    \text{Types } A &::=& \nat \alt \,\dyn \alt A \ra A' \alt A \times A'\\
    \text{Type Precision } c &::=& r(A) \alt c c' \alt \iarr \alt \inat \alt \itimes \alt c \ra c' \alt c \times c'\\
    \text{Values } V &::=& x \alt \upc c V \alt \zro \alt \suc\, V \alt \lda{x}{M} \alt (V,V') \\ 
    \text{Terms } M,N &::=& \err \alt \upc c M \alt \dnc c M \alt \zro \alt \suc\, M \alt \lda{x}{M} \\ 
     &&\alt M\, N \alt (M,N) \alt \textrm{let } (x,y) = M \textrm{ in } N \alt \synop{op}{V}{x}{M} \\
    \text{Contexts } \Gamma &::= &\cdot \alt \Gamma, x : A \\
    \text{Ctx Precision } \Delta &::=& \cdot\alt \Delta,x:c
  \end{array}

  \inferrule*[right=Up]
  {\Gamma \vdash M : A \and c : A \ltdyn A'}
  {\Gamma \vdash \upc c M : A'}

  \inferrule*[right=Dn]
  {\Gamma \vdash N : A' \and c : A \ltdyn A'}
  {\Gamma \vdash \dnc c N : A}

  \inferrule*[right=Eff]
  {\eff{op}{\sigma}{\tau} 
    \and \Gamma \vdash V : \sigma 
    \and \Gamma,x : \tau \vdash M : A}
  {\Gamma \vdash \synop{op}{V}{x}{M} : A}

  \inferrule{}{\Gamma \vdash \mho : A}
  \end{mathpar}
  \begin{mathpar}
    \inferrule{}{r(A) : A \ltdyn A}\and
    \inferrule{c : A \ltdyn A' \and c' : A' \ltdyn A''}{cc' : A \ltdyn A''}\and
    \inferrule{}{\iarr \colon \dyn \ra \dyn \ltdyn \dyn}\and
    \inferrule{}{\inat \colon \nat \ltdyn \dyn}\and
    \inferrule{}{\itimes \colon \dyn \times \dyn \ltdyn \dyn}\and
    \inferrule{c_i : A_i \ltdyn A_i' \and c_o : A_o \ltdyn A_o'}{c_i \ra c_o : (A_i \ra A_o) \ltdyn (A_i' \ra A_o')}\and
    \inferrule{c_1 : A_1 \ltdyn A_1' \and c_2 : A_2 \ltdyn A_2'}{c_1 \times c_2 : (A_1 \times A_2) \ltdyn (A_1' \times A_2')}\and
     r(A)c \equiv c\and
     c \equiv cr(A')\and
     c(c'c'') \equiv (cc')c''\and
     r(A_i \ra A_o) \equiv r(A_i) \ra r(A_o)\and
     r(A_1\times A_2) \equiv r(A_1) \times r(A_2)\and
     (c_i \ra c_o)(c_i' \ra c_o')\equiv (c_ic_i' \ra c_oc_o') \and
     (c_1\times c_2)(c_1'\times c_2')\equiv (c_1c_1' \times c_2c_2')
  \end{mathpar}
  \caption{Syntax for Gradual Typing with Effects}
  \label{fig:gtlc-syntax}
\end{figure}


% New, Licata, and Ahmed gave a syntax and logic axiomatizing the graduality
% property \cite{new-licata-ahmed2019}.

\section{Operational Semantics}\label{sec:operational-semantics}

In this section, we define a small-step operational semantics for the generic
cast calculus introduced in Section \ref{sec:syntax}. Our approach is based on
the recent work of Kavvos \cite{TODO}, who gave a generic operational semantics
for a language with algebraic effects.

Following the work of Kavvos, our operational semantics has two kinds of
reductions. There are \emph{pure reductions}, which comprise the usual $\beta$
reductions in an ordinary lambda calculus as well as reductions involving casts.
There are also \emph{impure reductions}, which result from the program making a
request to its environment and receiving a value in response.

The pure reductions are labelled by a natural number $i \in \{0, 1\}$, which
intuitively tracks the number of recursive unfoldings involved in each step. The
source of such steps is the downcast from the dynamic type to a function, since
in the denotational semantics the function is only available later and hence the
downcast involves an observable computational step. We track the recursive
unfolding of terms in the operational semantics to help us relate the
operational and denotational semantics, which we address in a subsequent
section.
%

The impure reductions do not involve any recursive unfolding, so they are not
labeled by a natural number. A general impure reduction has the form
%
\[ M \stepseff{op}{V}{W} N, \]
%
which means that $M$ sends value $V$ to the environment and continues as $N$
upon receiving response $W$.

% TODO if this is for fine-grain call-by-value, then do we need rules involving
% function application?
% TODO define evaluation contexts
% TODO need to define a subset of "first-order values" to be used in the effect
% operations.

\begin{figure}

  \textbf{Beta rules:}

  \begin{mathpar}

    % Usual beta rules:
    \inferrule*[right=Beta]
      { }
      {(\lambda x : A.M)(V) \stepsin{0} M[V/x]}

    \inferrule*[right=Split]
      { }
      {\synsplit{(V_1, V_2)}{x_1}{x_2}{M} \stepsin{0} M[V_1, V_2 / x_1, x_2]}

    \inferrule*[right=Bind]
      { }
      {\synbind{\ret\, V}{x}{M} \stepsin{0} M[V/x]}
  \end{mathpar}

  \textbf{Casts (functorial actions and composition):}
  \begin{mathpar}

    % Upcast of a function steps when applied to a value:
    \inferrule*[right=Up-Arr]
      { c_i : A_i \ltdyn A_i' \and c_o : A_o \ltdyn A_o' 
        \and \Gamma \vdash V_f : A_i \ra A_o 
        \and \Gamma \vdash V : A_i'}
      { ((\upc (c_i \ra c_o) V_f)\, V) \stepsin{0} (\upc {c_o} (V_f\, (\dnc {c_i} V))) }

    % Downcast of a function steps when applied to a value:


    % Up/down cast of a paired value steps to the pair of the up/down casted
    % values:
    \inferrule*[right=Up-Pair]
      { c_1 : A_1 \ltdyn A_1' \and c_2 : A_2 \ltdyn A_2'}
      { \upc (c_1 \times c_2) (V_1, V_2) \stepsin{0} (\upc {c_1} V_1, \upc {c_2} V_2) }

    \inferrule*[right=Dn-Pair]
      { c_1 : A_1 \ltdyn A_1' \and c_2 : A_2 \ltdyn A_2'}
      { \dnc (c_1 \times c_2) (V_1, V_2) \stepsin{0} (\dnc {c_1} V_1, \dnc {c_2} V_2) }
  

    % "Composition" of casts with inj-arr and inj-times:
    \inferrule*[right=$\iarr$CompUp]
      { c : A \ltdyn (\dyntodyn) \and \Gamma \vdash M : A }
      { (\upc (\injarr c) M) \stepsin{0} (\upc \iarr (\upc c M))}

    \inferrule*[right=$\iarr$CompDn]
      { c : A \ltdyn (\dyntodyn) \and \Gamma \vdash M : \dyn }
      { (\dnc (\injarr c) M) \stepsin{0} (\dnc c (\dnc \iarr M))}

    \inferrule*[right=$\itimes$CompUp]
      { c : A \ltdyn (\dyn \times \dyn) \and \Gamma \vdash M : A }
      { (\upc (\injtimes c) M) \stepsin{0} (\upc \itimes (\upc c M))}

    \inferrule*[right=$\itimes$CompDn]
      { c : A \ltdyn (\dyn \times \dyn) \and \Gamma \vdash M : \dyn }
      { (\dnc (\injtimes c) M) \stepsin{0} (\dnc c (\dnc \itimes M))}



  \end{mathpar}

  \textbf{Casts involving dyn:}
  \begin{mathpar}
    % Downcasts of upcasted-to-dyn values:
    \inferrule*
      { }
      { \dnc \inat (\upc \inat V) \stepsin{0} V}

    \inferrule*
      { c : A \ltdyn (\dyn \times \dyn)}
      { \dnc {\injtimes c} (\upc {\injtimes c} V) \stepsin{0} V}

    \inferrule*
      { c : A \ltdyn (\dyntodyn)}
      { \dnc {\injarr c} (\upc {\injarr c} V) \stepsin{1} V}

    \inferrule*
     { \text{tag}_1, \text{tag}_2 \in \{ \mathbb{N}, \times, \ra \} }
     { \dnc {\text{Inj}_{\text{tag}_1}} (\upc {\text{Inj}_{\text{tag}_2}} V) \stepsin{0} \mho }


  \end{mathpar}

  \textbf{Effects and congruence rules:}
  \begin{mathpar}

    % Effects
    \inferrule*[right=Eff]
      { \eff{op} \sigma \tau }
      { \synop{op} V x M \stepslabel{\eff {op} V W} M[W/x] }

    % Congruences
    \inferrule*[right=Cong]
      { i \in \{0,1\} \and M \stepsin i N}
      { E[M] \stepsin{i} E[N] }

    \inferrule*[right=EffCong]
      {M \stepslabel {\eff {op} V W} N}
      {E[M] \stepslabel {\eff {op} V W} E[N]}


  \end{mathpar}
  \caption{Operational Semantics (Selected Rules)}
  \label{fig:operational-semantics}
\end{figure}

There are $\beta$ rules for functions, pairs, and bind, which are standard. Next
are the rules involving casts. We first have rules involving upcasts and
downcasts of functions and products.
%
When a function $V_f$ is upcasted, it doesn't immediately step. When the
upcasted function is applied to a value, the value is downcasted according to
the domain and applied to the original function $V_f$, and then the result is
upcasted according to the codomain. The dual holds for a function that is
downcasted.
%
When a pair of values $(V_1, V_2)$ is upcasted, we simply upcast each component
of the pair, and likewise for downcasting.
%
Next we have rules for up/down casts involving composed precision derivations,
i.e., $\injarr c$ and $\injtimes c$. These rules decompose such casts into the
cast for $c$ and the primitive injection casts $\iarr$ or $\itimes$. 
% These rules express the behavior of such casts in
% terms of the behavior of the cast for $c$ and the behavior of the primitive
% injection casts $\iarr$ or $\itimes$. 
For example, to upcast a term using the precision derivation $\injarr c$ where
$c : A \ltdyn (\dyntodyn)$, we first upcast the term by $c$, and then upcast the
result by $\iarr : \dyntodyn \ltdyn \dyn$.
%
Finally, we have the rules involving downcasting a value that was cast to $\dyn$
using the injection casts $\inat$, $\injarr c$, and $\injtimes c$. In each case,
we check that the tag of the upcast matches the tag of the downcast, e.g., the
upcast is $\iarr$ and the downcast is $\iarr$. If the tags match, then the cast
steps to the original value; otherwise, the result is an error. Note that when
downcasting a dynamically-typed function, we record one recursive unfolding. In
the denotational semantics, this corresponds to inserting a $\theta$.
%
% Note also that these rules only mention the primitive injection forms $\inat$,
% $\iarr$, and $\itimes$ rather than the composed forms $\injarr c$ and $\injtimes
% c$. This is sufficient because casts involving $\injarr c$ or $\injtimes c$ will
% first step to terms involving $\iarr$ and $\itimes$, which can then be handled
% using the rules just discussed.

Next, we have our first impure reduction rule. This rule states that for an
operation $\text{op}$ with request type $\sigma$ and response type $\tau$, then
for any possible request value $V$ and response value $W$, the term $\synop{op}
V x M$ steps to $M[W/x]$. This represents the fact that the environment has
responded with value $W$, after which the program continues as $M$ with $W$ in
place of $x$.
%
The final two rules are simply congruence rules allowing us to perform pure and
impure reductions under an arbitrary evaluation context.

\subsection{A Guarded Interaction Tree Semantics}


\section{Denotational Semantics}\label{sec:denotational-semantics}

In this section, we discuss how to adapt the model of Giovannini et al.
\cite{gtt-sgdt} to accommodate algebraic effects. Recall that in their model,
pure values were modeled as elements of a \emph{predomain}, and effectful
computations were modeled as elements of an \emph{error domain}, which is a
predomain equipped with algebraic operations supporting error and computational
steps. Our approach will be to generalize the notion of error domain to be an
algebra of a suitable endofunctor corresponding to the signature of effects
being modeled. Then the free error domain construction will be generalized to
construct the free such algebraic structure, as an \emph{inductive-guarded
interaction tree}.

In addition to generalizing error domains, we need to make an important
technical modification to the model of Giovannini et al. This change is
motivated by our goal of relating the denotational and operational semantics via
an adequacy result. In op. cit., the authors mention that the denotations of
casts in their model involves additional perturbations (i.e., computational
steps) in order to place certain terms involving casts in lock-step. These
additional steps ultimately did not matter for the final result, since the
graduality property is step-insensitive. However, intensionally, these casts
behave differently from the corresponding casts without the extra steps.
%
This presents a difficulty when relating their denotational model to an
operational semantics: in order for the two semantics to be in lock-step, we
must define the operational semantics in such a way that the casts take
additional steps just as they do in the denotational semantics. But this
requirement is not desirable, as it restricts the kinds of operational semantics
to which our theorem applies to those that happen to feature this strange cast
semantics. Users of our framework would need to contort their operational
semantics into this form if they would like to apply the adequacy result.

Our remedy for this situation is to note that there is an additional degree of
freedom in the model of Giovannini et al. that up to this point has remained
unexploited. Leveraging this, we define two \emph{different} cast semantics, one
where the casts involve perturbations, and one where they do not. The former is
used in constructing the semantic proofs of the cast axioms, where being in
lock-step is crucial. The latter is what we use for the denotational term
semantics of casts, allowing the denotations of casts to be in lock-step with
the casts in a typical operational model, i.e., where the casts do not take
additional steps. While the intensional behavior of these two kinds of casts are
distinct, they are extensionally related in that they are \emph{weakly
bisimilar}, which is sufficient for the final proof of graduality.

\subsection{Generalizing Error Domains}

In this section, we discuss how to build a model of gradual typing from scratch
given a collection of effects presented algebraically. Our goal will be to
generalize the notion of \emph{error domain} from prior work to accommodate
these effects.
%
We begin by reviewing the key ingredients in the denotational semantics of
Giovannini et al. \cite{gtt-sgdt}. For further discussion, see their paper.

\begin{definition}
A \textbf{predomain} $A$ consists of a set $A$ along with two relations:
\begin{itemize}
    \item A partial order $\ltdyn_A$.
    \item A reflexive, symmetric ``bisimilarity'' relation $\bisim_A$.
\end{itemize}
\end{definition}

\begin{definition}
An \textbf{error domain} $B$ consists of a predomain $B$ along with the
following data:
\begin{itemize}
    \item A distinguished ``error'' element $\mho_B \in B$ such that $\mho
    \ltdyn_B x$ for all $x$
    \item A morphism of predomains $\theta_B \colon \later B \to B$
    \item For all $x, y : B$ with $x \bisim_B y$, we have $\theta_B(\nxt\, x)
    \bisim_B y$
\end{itemize}
\end{definition}

A \textbf{morphism of predomains} $f : A \to A'$ is a function between the
underlying sets that is both \emph{monotone} (if $x \ltdyn_A x'$, then $f(x)
\ltdyn_{A'} f(x')$), and \emph{bisimilarity-preserving} (if $x \bisim_A x'$,
then $f(x) \bisim_{A'} f(x')$).

A \textbf{morphism of error domains} is a morphism of the underlying predomains
that preserves the error element and the $\theta$ map.
%
For an error domain $B$, we define the predomain morphism $\delta_B := \theta_B
\circ \nxt$.

A \emph{predomain relation} $c : A \rel A'$ is a relation between $A$ and $A'$
that is upward closed in $A'$ and downward closed in $A$, as defined in the
previous section (no condition is required for bisimilarity). The ``reflexive''
predomain relation is written $r(A)$ and is given by the ordering $\ltdyn_A$.
The relational composition of two predomain relations is also a predomain
relation, and $r(A)$ is the identity for this composition.

Error domain relations are predomain relations that are additionally a
congruence with respect to the algebraic operations $\theta$ and $\mho$:
%
\begin{definition}
  An \textbf{error domain relation} $d : B \rel B'$ is a relation $d$ on the
  underlying predomains satisfying
  \begin{enumerate}
     \item ($\mho$ congruence): $\mho_B \binrel{d} \mho_{B'}$.
     \item ($\theta$ congruence): For all $\tilde{x}$ in $\laterhs B$ and
     $\tilde{y} \in \laterhs B'$, if $\laterhs_t (\tilde{x}_t \binrel{d}
     \tilde{y}_t)$, then $(\theta_B (\tilde{x}) \binrel{d} \theta_{B'}
     (\tilde{y}))$. 
  \end{enumerate}
\end{definition}
%
There is again a reflexive error domain relation $r(B)$ given by the ordering.
The relational composition of two error domain relations in general is not an
error domain relation ($\theta$ congruence in particular is not preserved), but
we instead use an alternative ``free composition'' $dd' : B \rel B''$ that is
inductively generated from the relational composition.

Lastly, we have the notion of a square $f \ltsq{c_i}{c_o} g$ between predomain
morphisms and of a square between error doamin morphisms:

\begin{definition}
Let $f : A_i \to A_o$ and $g : A_i' \to A_o'$ be predomain morphisms, and let
$c_i : A_i \rel A_i'$ and $c_o : A_o \rel A_o'$ be predomain relations. There is
a \emph{predomain square} $f \ltsq{c_i}{c_o} g$ if for all $x \binrel{c_i} y$,
we have $f(x) \binrel{c_o} g(y)$. 

Likewise, an \emph{error domain square} $\phi \ltsq{d_i}{d_o} \psi$ is defined
similarly, with no extra interaction with the algebraic structure.
\end{definition}

We next review how the model of Giovannini et al. validates the graduality
property. The authors of that work took as a starting point a semantics with a
notion of representable relations, which in prior work was sufficient for
interpreting the four cast rules crucial for graduality. However, as the authors
show, the notion of representability cannot be applied unchanged to the guarded
setting because of a fundamental tension between transitivity and
step-insensitive reasoning. As a result, Giovannini et al. worked with a
weakened notion that they called \emph{quasi-representability}, which is
essentially representability up to an explicit delay term. They then developed a
theory of these delay terms, which they called \emph{syntactic perturbations}.

% TODO more detail on the definition of quasi-representability?


In the remainder of this section, we discuss the changes we must make to
accommodate arbitrary algebraic effects. It turns out that they key question is
how to generalize the notion of error domain.

Consider a signature $\Sigma$ of algebraic effects given by a type $\Op$ of
operation names, and for each $\oper : \Op$ a pair of \emph{predomains}
$\opin{\oper}$ and $\opout{\oper}$ representing respectively the request and
response types for $\oper$. This signature induces a polynomial endofunctor
$\Sigma$ on the category of predomains and predomain morphisms.
%
Our generalization of error domains will be an algebra $B$ of the endofunctor
$\Sigma$, equipped with a distinguished error element $\mho$ and a distinguished
$\theta_B \colon \later B \to B$. 
%
Since $\Sigma$ acts on predomains, an algebra of the endofunctor is a predomain
$B$ equipped with a morphism of predomains $\Sigma(B) \to B$. This means in
particular that $B$ has ordering and bisimilarity relations.
%
As in the definition of an error domain from prior work, we require that $\mho$
is the smallest element in the ordering, that $\theta$ respects the relations
(i.e., it is a morphism of predomains), and that $\delta_B := \theta_B \circ
\nxt$ is bisimilar to the identity.

% As in the definition of an error domain, we also require that the algebra comes
% with ordering and a bisimilarity relations, that $\mho$ is the smallest element
% in the ordering, that $\theta$ respects the relations (i.e., it is a morphism of
% predomains), and that $\delta_B := \theta_B \circ \nxt$ is bisimilar to the
% identity.


% To model gradual typing, we must ensure that the signature has cases for errors
% and computational steps. To this end, given a signature $\Sigma$ as above, we
% define $\Sigmahat := (\hat{S}, \hat{P})$, where $\hat{S} = S + \bool$, and
% $\hat{P}$ is given by $\hat{P}(\inl s) = P\, s$, $\hat{P}(\inr \text{true}) =
% \bot$, and $\hat{P}(\inr \text{false}) = \top$. 
% %
% We will then be interested in algebras of $\Sigmahat \circ \later$. Here, we
% note that this $\later$ is operating on predomains. Furthermore, this definition
% implies that all occurrences of the input $X$ are guarded by a $\later$. A
% further generalization would be to consider algebras of a functor where some
% occurrences of $X$ are unguarded and others are guarded. In this case, we would
% need to assume more about our given functor $\Sigma$, e.g., that it is presented
% as a \emph{polynomial} endofunctor; we will pursue this in the subsequent
% section.

% More concretely, an algebra of the functor $\Sigmahat \circ \later$ consists of
% a predomain $|B|$ along with a morphism of predomains
% %
% $\alg_B \colon \Sigmahat (\later |B|) \to |B|$.
% %
% Note that $\Sigmahat (\later |B|) = \Sigma(\later |B|) \uplus 1 \uplus (\later
% |B|)$. Thus, an algebra map $\alg_B$ induces three morphisms of predomains; we
% denote these by $\oper_B \colon \Sigma(\later |B|) \to |B|$, $\mho_B \colon 1
% \to |B|$, and $\theta_B \colon \later |B| \to |B|$.
% %
% Notice also that, by definition of morphism of predomains, these maps preserve
% the ordering and bisimilarity relations of the relevant predomains.

% TODO: update this

More formally, we make the following definition:
%
\begin{definition}\label{def:sigma-domain} 
  
  Let $\Sigma = (\Op, \text{In}, \text{Out})$ be a signature as defined above. A
  \emph{$\Sigma$-domain} is a predomain $B$ along with, for each $\text{op} \in
  \Op$, a morphism of predomains
  
  \[ \alpha^{\oper} : \opin{\oper} \times (\opout{\oper} \to B) \to B. \]
  
  Furthermore, the predomain $B$ has a distinguished element $\mho_B$ such that
  $\mho_B \ltdyn_B x$ for all $x \in B$, along with a distinguished morphism of
  predomains $\theta_B : \later B \to B$, such that if $x \bisim_{B} y$, then
  $\delta_B \bisim_{B} y$ where $\delta_B := \theta_B \circ \nxt$.
\end{definition}

As with ordinary error domains, there are corresponding notions of morphisms,
relations, and squares of $\Sigma$-domains. Before introducing these, we first
discuss the action of $\Sigma$ on predomain morphisms and predomain relations.

Let $c : A \rel A'$ be a predomain relation. The action of the endofunctor
$\Sigma$ on $c$, written $\Sigma(c)$ is a predomain relation between $\Sigma(A)$
and $\Sigma(A')$.



Now we can introduce the definitions of morphisms and relations of
$\Sigma$-domains.
%
\begin{definition}\label{def:sigma-domain-mor-rel-sq}
  Fix $\Sigma$ as above, and let $B$ and $B'$ be $\Sigma$-domains. A
  \emph{morphism of $\Sigma$-domains} from $B$ to $B$' is a morphism $\phi$ of
  the underlying predomains that preserves the algebraic structure.
  %
  Spelling this out concretely, we have:
  %
  (1) $\phi (\alpha^{\oper}_B(x, k)) = \alpha^{\oper}_{B'}(x, \phi \circ k)$,
  %
  (2) $\phi(\mho_B) = \mho_{B'}$, and
  %
  (3) $\phi \circ \theta_B = \theta_{B'} \circ (\later \phi)$.

  A \emph{relation of $\Sigma$-domains} $d : B \rel B'$ is a relation $|d|$ of
  the underlying predomains that is a congruence with respect to the algebraic
  structure. In particular:
  %
  (1) If $x \binrel{(\Sigma(|d|))} y$, then $\alpha^{\oper}_B(x) \binrel{|d|} \alpha^{\oper}_{B'}(y)$.
  %
  (2) $\mho_B \binrel{|d|} \mho_{B'}$, and
  %
  (3) If $\later_t(\tilde{x}_t \binrel{|d|} \tilde{y}_t)$, then
  $\theta_B(\tilde{x}) \binrel{|d|} \theta_{B'}(\tilde{y})$

  A \emph{square of $\Sigma$-domains} $\phi \ltsq{d}{d'} \phi'$ is defined
  as a square of the underlying predomains. 
\end{definition}

\begin{remark}
  The requirement that $\mho_B$ be related to $\mho_{B'}$ in the definition of
  relation may seem weaker than the corresponding requirement in the definition
  of a relation of error domains in the previous section, where we required that
  $\mho_B \binrel{|d|} y$ for \emph{all} $y \in B'$. However, the former implies
  the latter using the fact that predomain relations are closed under the
  ordering on both sides and that $\mho_{B'} \ltdyn y$ for all $y$.
\end{remark}





\section{Relating the Two Semantics}\label{sec:relating-the-semantics}

Having specified the operational and denotational semantics, we now seek to
prove a theorem relating them. Our goal is to show that the denotatinoal model
is \emph{adequate} for the operational semantics, that is, if the denotations of
two terms are equal, then their operational behavior is the same.

To this end, we will define a logical relation between (closed) syntactic terms
of type $A$ and their denotations as elements of $F\sem{A}$, where $F$ is the
functor constructing the free error-effect algebra as discussed in Section
\ref{sec:denotational-semantics}.

\section{Discussion}\label{sec:discussion}




\bibliographystyle{ACM-Reference-Format}
\begin{DIFnomarkup}
\bibliography{\bibliographypath}
\end{DIFnomarkup}

% EXTENDED VERSION
\newpage
\appendix
% \input{appendix}

\end{document}
